\section{Stationary equilibrium}\label{sec:eqm}

Every cohort is born with no assets and an earnings state drawn from a fixed law $\nu$, lives $J$
periods, and is replaced; cohorts have equal weight. The distribution of the youngest cohort is
therefore the same at every interest rate, and the distribution of every older cohort is that law
pushed forward by the policies of the ages in between. This is the structural advantage of the
life cycle over the infinite horizon, and it removes most of the machinery the stationary case
needed.

\subsection{The aggregate is a finite composition}

Because assets at birth are degenerate at zero, no measure theory is required. Let
\[
  (\mathcal{S}_k h)(a,z)=\sum_{z'}\pi(z,z')\,h\bigl(g_k(a,z),z'\bigr)
\]
be the expectation of $h$ next period for a household with $k$ periods still to come
(\lean{stageStep}). The assets a household holds over the rest of its life, summed across ages
and seen from its current state, satisfy
\[
  A_0(a,z)=a,\qquad A_{k+1}(a,z)=a+(\mathcal{S}_{k+1}A_k)(a,z)
\]
(\lean{cohortAssets}), and aggregate capital per head is
$K=\frac1J\sum_z\nu(z)\,A_{J-1}(0,z)$ (\lean{olgCapital}). There is no pushforward measure, no
Feller property, no stationary distribution, no Doeblin condition and no uniqueness of an
invariant law: the equilibrium object is a finite sum of finite sums, and the recursion above is
its definition.

\subsection{The chain}

\begin{theorem}[Light's Theorem 2 for the life cycle]\label{thm:thm2}
Suppose two economies share the earnings chain and their policies are ordered at every age,
$g^{(1)}_k\le g^{(2)}_k$ on the asset region. Then $A^{(1)}_k\le A^{(2)}_k$ for every $k$, and
aggregate capital is ordered, $K^{(1)}\le K^{(2)}$.
\formal{IncomeFluctuation.cohortAssets\_le,
IncomeFluctuation.olgCapital\_le\_of\_stagePolicy\_le}
\end{theorem}

The proof is an induction on age with two ingredients, and both are already available. The class
of functions increasing in assets is preserved by $\mathcal{S}_k$
(\lean{monoRegion\_stageStep}), because next period's assets rise with this period's at each
earnings state and for any continuation, which is Theorem~\ref{thm:azmono}. And ordered policies
keep ordered functions ordered (\lean{stageStep\_le\_stageStep}): the higher-rate economy
evaluates a larger function at a larger asset level. Note what is \emph{not} needed. The
infinite-horizon version of this step required monotone transitions, a coupling argument and
first-order stochastic dominance of the stationary distributions; here the whole statement is an
inequality between two functions on the same finite state space, because the distribution the two
economies start from is literally the same.

\subsection{A floor on saving from primitives}

Existence needs capital supply at the top of the rate interval to exceed demand. In the infinite
horizon the only floor available came from an ergodic identity that loses a factor
$\|u'\|_\infty^2/\operatorname{Var}$, with the consequences for the constants reported in the
companion write-up. The life cycle needs nothing of the kind: the upper sandwich of
Theorem~\ref{thm:sandwich} bounds consumption at every age, and what is not consumed is saved.

\begin{proposition}[Affine floor on saving]\label{prop:floor}
At every age and every state in the region,
\[
  g_k(a,z)\ \ge\ (1-\kappa_k)\bigl(y_{\min}+Ra\bigr)-\kappa_k H_k(z).
\]
\formal{IncomeFluctuation.le\_stagePolicy}
\end{proposition}

Iterated from zero assets at birth this is an explicit floor on the assets of every cohort, in
the primitives alone: the discount factor, the rate, the earnings bounds and the horizon.

\subsection{Uniqueness}

\begin{theorem}[Single crossing]\label{thm:thm3}
A capital supply that never falls with the rate meets a strictly falling demand at most once.
\formal{IncomeFluctuation.olgEquilibriumRate\_unique}
\end{theorem}

With Theorem~\ref{thm:thm2} this reduces uniqueness of the stationary equilibrium to Light's
Theorem 1, that saving rises with the rate at every age, which is the subject of
Section~\ref{sec:thm1}: it holds at Aiyagari's calibration for relative risk aversion up to about
$3.4$, so the chain closes at his $\mu=3$ and not at his $\mu=5$.

What is still missing for existence is continuity of capital supply in the rate. In the
infinite horizon that was a substantial development, because the stationary distribution had to
be shown to move continuously; here it is continuity of finitely many policies in the rate,
composed. That is the next piece.
